\chapter{El dolor: una señal de alarma}
\chaptermark{El dolor}
\label{sec:pain}
\begin{chaptergoals}
\item por qué el dolor es una señal de alarma de alta prioridad y no un canal de medición;
\item cómo un cable rápido y uno lento informan de dónde está el daño y de qué tan grave es;
\item cómo la compuerta medular deja al tacto vetar el dolor, y la difusión fija el volumen de alarma;
\item cómo la sensibilización sube la ganancia tras una lesión, y cómo el dolor enseña e interrumpe.
\end{chaptergoals}

\section{Alarma, no medición}

Todos los sensores del capítulo~\ref{sec:sensors} medían el mundo: cuánta luz, qué tono, qué presión. El dolor funciona de otra manera. No informa “qué hay ahí”, sino “peligro, déjalo todo”. Para un ingeniero no es un canal de medición, sino una \emph{señal de alarma}: una línea de máxima prioridad que debe abrirse paso a través de cualquier trabajo en curso de la máquina.

Un sistema de alarma se distingue de un canal de medición por tres propiedades, y el dolor tiene las tres. Primera: es difícil atenuarlo por adaptación. Los sensores del dolor se adaptan mucho menos que los del tacto, y tras una lesión leve incluso se vuelven más sensibles~\cite{LaMotte1982hyper}; el debilitamiento de la respuesta tras un calentamiento intenso se ha medido solo en uno de sus tipos~\cite{Treede1995heat}. El sensor de luz se calla ante un fondo constante (fig.~\ref{fig:adaptation}), pero una alarma que se calla al minuto no sirve. Segunda: tiene un umbral alto. Los sensores del dolor se activan solo cuando el estímulo se acerca a lo peligroso; por ejemplo, ante el calor, los umbrales van de $41^\circ$ a $46$--$48^\circ$ según el tipo de sensor~\cite{Treede1995heat, Weidner1999cfib}. Tercera, y la principal para este capítulo: el volumen de la alarma no lo decide solo el sensor, sino también la propia máquina. Una misma herida duele distinto en distintas circunstancias. Veamos con qué piezas ya conocidas está armado esto.

Los sensores del dolor son neuronas \nt{GLUT}, como las demás neuronas sensoriales del capítulo~\ref{sec:sensors}. Algunas de ellas, junto con el glutamato, liberan la sustancia~P del apéndice~\ref{sec:othermediators}, que refuerza lentamente la transmisión.

\section{Dos cables: dolor rápido y dolor lento}

Golpéese un dedo y sentirá el dolor dos veces: primero uno breve y agudo, y un segundo después otro sordo y prolongado. Son dos cables distintos. El cable rápido está aislado y conduce las espigas a una velocidad $v$ de $5$ a $30$ metros por segundo; el lento es fino y desnudo, $0.5$--$2$ metros por segundo. La señal del pie recorre alrededor de un metro, y el tiempo de llegada
\begin{empheq}[box=\keybox]{equation}
t \;=\; \frac{L}{v}
\label{eq:paintime}
\end{empheq}
es de decenas de milisegundos para el cable rápido y de alrededor de un segundo para el lento. Los dos cables llevan dos mensajes. El rápido dice \emph{dónde} y \emph{cuándo}: sus espigas llegan antes y con más precisión, como en el código temporal del capítulo~\ref{sec:code}. El lento dice \emph{qué tan grave}, y su dolor sordo y prolongado mantiene a la máquina en modo de alarma mientras la lesión no haya pasado.

\section{La compuerta en la médula espinal}

El primer lugar al que llega la señal de dolor es la médula espinal, y allí pasa por una \emph{compuerta}~\cite{MelzackWall1965}; las neuronas concretas de esta compuerta se encontraron hace relativamente poco~\cite{Duan2014}. En la neurona \nt{GLUT} de salida, que transmite el dolor hacia el cerebro, convergen el cable del dolor y el cable del tacto. Y al lado hay una neurona intermediaria inhibidora, \nt{GABA} o \nt{GLYC}, a la que excita el tacto (fig.~\ref{fig:gate}). El tacto cierra la compuerta. En la teoría original de la compuerta~\cite{MelzackWall1965}, el dolor, al contrario, apaga esa intermediaria, y un dolor fuerte abre la compuerta; es un supuesto de la teoría que no se ha mostrado directamente en las neuronas de la compuerta encontradas.

\begin{figure}[t]
\centering
\begin{tikzpicture}[>=Stealth,
  nrn/.style={circle,draw,very thick,minimum size=1.0cm,inner sep=0pt,font=\scriptsize},
  inh/.style={-{Bar[width=7pt]},thick,red!70!black},
  bc/.style={->,thick,densely dashed,orange!85!black}]
\node[nrn,fill=blue!6] (Z) at (6,0) {\nt{GLUT}};
\node[nrn,fill=red!10] (Q) at (3.6,-1.6) {\nt{GABA}};
\draw[->,very thick,red!70!black] (0,0.6) node[left,font=\small,black,align=right] {dolor $\nu$} -- (5.45,0.15);
\draw[->,very thick,blue!60!black] (0,-2.6) node[left,font=\small,black,align=right] {tacto $\mu$} -- (2.9,-2.0);
\draw[->,thick,blue!60!black] (1.8,-2.35) -- (5.5,-0.35) node[pos=0.8,below right=-2pt,font=\scriptsize] {$w_{z\mu}$};
\draw[inh] (1.6,0.5) -- (3.35,-1.1) node[pos=0.5,left=1pt,font=\scriptsize] {$w_{q\nu}$};
\draw[inh] (Q) -- (Z) node[pos=0.5,above left=-1pt,font=\scriptsize] {$\beta$};
\draw[->,very thick] (Z) -- (8.2,0) node[right,font=\small,align=left] {al cerebro: $z$};
\draw[bc] (6,2.1) node[above,font=\scriptsize,black,align=center] {desde arriba: \nt{SERO}, \nt{NORA}, opioides} -- (Z.north) node[pos=0.45,right,font=\scriptsize,black] {ganancia $g$};
\end{tikzpicture}
\caption{La compuerta del dolor en la médula espinal según el modelo~\eqref{eq:gate}. La neurona \nt{GLUT} de salida recibe el dolor $\nu$ y una excitación débil del tacto $\mu$. La intermediaria inhibidora (\nt{GABA} o \nt{GLYC}) se excita con el tacto y, según el supuesto de la teoría de la compuerta, el dolor la apaga. Desde arriba, por los canales de difusión, llega la ganancia $g$.}
\label{fig:gate}
\end{figure}

Escribámoslo para las frecuencias, como en el capítulo~\ref{sec:gaba}. La frecuencia de la intermediaria $q$ y la frecuencia de salida $z$ son
\begin{empheq}[box=\keybox]{equation}
q \;=\; \bigl[w_{q\mu}\,\mu + q_0 - w_{q\nu}\,\nu\bigr]_+ ,
\qquad
z \;=\; \bigl[\,g\,(\nu + w_{z\mu}\,\mu) - \beta\,q\,\bigr]_+ ,
\label{eq:gate}
\end{empheq}
donde $\nu$ es la entrada del dolor, $\mu$ la entrada del tacto, $w$ los pesos de las conexiones (el primer índice indica a quién; el segundo, de quién), $\beta$ la fuerza de la inhibición, $q_0$ la actividad de fondo propia de la intermediaria y $g$ la ganancia que se fija desde arriba (sobre ella, más abajo). Es el veto~\eqref{eq:veto} del capítulo~\ref{sec:gaba}, “dolor, pero no con tacto”, con una corrección tomada de la teoría de la compuerta como supuesto: el dolor fuerte apaga por sí mismo a la neurona que veta.

Calculamos el modelo con $w_{q\mu}=1$, $q_0=0.2$, $w_{q\nu}=0.5$, $w_{z\mu}=0.3$, $\beta=1.5$ (fig.~\ref{fig:gatecurves}). Sin tacto, el dolor empieza a pasar desde $\nu\approx0.18$. Con tacto, el umbral sube a $0.86$, y con $\nu=1$ la salida cae a la cuarta parte, de $1$ a $0.25$. Pero con un dolor fuerte, $\nu=2$, el tacto ya no cambia nada: la compuerta está abierta de par en par. Así pasa en la vida: frotar un golpe ayuda, pero ante una lesión seria, no. El tacto por sí solo, sin dolor, no pasa por la compuerta: la alarma no salta por un roce.

\begin{figure}[t]
\centering
\begin{tikzpicture}[>=Stealth,x=5.2cm,y=1.35cm]
\draw[->,gray!70!black] (0,0) -- (2.12,0) node[right,font=\small,black] {dolor $\nu$};
\draw[->,gray!70!black] (0,0) -- (0,4.3) node[left,font=\small,black] {$z$};
\foreach \t in {0,0.5,1,1.5,2}{\draw[gray!60!black] (\t,-0.05) -- (\t,0.05); \node[below,font=\scriptsize] at (\t,0) {\t};}
\foreach \f in {1,2,3,4}{\draw[gray!60!black] (-0.01,\f) -- (0.01,\f); \node[left,font=\scriptsize] at (0,\f) {\f};}
\draw[very thick,red!70!black] plot file {figures/pain_gate_notouch.dat};
\draw[very thick,blue!60!black] plot file {figures/pain_gate_touch.dat};
\draw[very thick,orange!85!black,densely dashed] plot file {figures/pain_gate_sens.dat};
\draw[very thick,red!70!black] (0.08,3.9) -- (0.2,3.9) node[right,font=\scriptsize,black] {sin tacto};
\draw[very thick,blue!60!black] (0.08,3.45) -- (0.2,3.45) node[right,font=\scriptsize,black] {con tacto};
\draw[very thick,orange!85!black,densely dashed] (0.08,3.0) -- (0.2,3.0) node[right,font=\scriptsize,black] {tras la sensibilización, $g=2$};
\end{tikzpicture}
\caption{Salida de la compuerta~\eqref{eq:gate} en función de la intensidad del dolor. El tacto sube el umbral y atenúa el dolor débil y el moderado, pero no el fuerte. La sensibilización (línea discontinua) duplica la ganancia: el mismo dolor se transmite el doble de fuerte, y el umbral baja.}
\label{fig:gatecurves}
\end{figure}

\section{La perilla de volumen desde arriba}

La compuerta no se controla solo desde abajo. Desde el tronco encefálico le llegan vías descendentes que cambian la ganancia $g$ en~\eqref{eq:gate}: \nt{SERO}, \nt{NORA} y los péptidos opioides del apéndice~\ref{sec:othermediators}~\cite{Fields2004}. Son los mismos canales de difusión de los capítulos~\ref{sec:serocomputer}--\ref{sec:acet}, solo que aquí su perilla es el volumen de la alarma. Pueden girarla en los dos sentidos: los mismos circuitos descendentes pueden tanto atenuar el dolor como amplificarlo.

¿Para qué querría la máquina bajar la alarma? Porque la alarma es una interrupción, y una interrupción no siempre es oportuna. Un animal que huye de un depredador no debe detenerse por una pata herida: primero correr, luego lamerse la herida. La expectativa de alivio también atenúa el dolor, y parte de ese efecto desaparece si se bloquean los receptores opioides~\cite{Levine1978}: la expectativa, calculada en el cerebro, baja por el canal opioide hasta la compuerta. Este cuadro en sí está medido; cómo decide exactamente el cerebro cuál debe ser el volumen es una hipótesis.

\section{Sensibilización: una alarma que aprende}

Tras una lesión, las neuronas de salida de la compuerta se vuelven más sensibles: la misma entrada provoca una salida mayor, y el umbral baja~\cite{Woolf1983}. En el modelo~\eqref{eq:gate} esto es un aumento de la ganancia $g$, y su descripción más simple es un circuito RC lento que carga el propio dolor:
\begin{empheq}[box=\keybox]{equation}
\tau_s\,\frac{\dd g}{\dd t} \;=\; -(g-1) + k_s\,z ,
\label{eq:sensitization}
\end{empheq}
donde $\tau_s$ es el tiempo en que la sensibilidad vuelve a la normalidad, más largo que la duración del propio dolor, porque la sensibilización persiste incluso después de que cesó el estímulo lesivo~\cite{Woolf1983}, y $k_s$ es la fuerza de la carga. Mientras el tejido está lesionado, la salida de la compuerta mantiene $g$ por encima de uno, y todo lo que ocurre cerca de la herida se transmite más fuerte (línea discontinua de la fig.~\ref{fig:gatecurves}: con $g=2$ el umbral baja a $0.11$ y la salida se duplica). Es un ajuste razonable de la alarma: mientras la lesión no sane, mejor pecar de precavido.

Para un ingeniero, es una realimentación positiva con una fuga lenta: el dolor sube la ganancia, y la ganancia, el dolor. Mientras el lazo es débil, $g$ vuelve a la normalidad cuando la herida sana. Pero si el lazo es fuerte, la alarma puede quedarse trabada en encendido incluso cuando la lesión ya no existe, como el grupo con realimentación fuerte de la fig.~\ref{fig:bistable}. El modelo~\eqref{eq:sensitization} es una simplificación; que la sensibilización central existe está medido.

\section{El dolor como maestro y como interrupción}

En la máquina, el dolor tiene dos trabajos además de la alarma.

\emph{Maestro.} El dolor es una recompensa negativa: en la fórmula del error~\eqref{eq:ctd} entra como $r<0$. Todo lo dicho en el capítulo~\ref{sec:dofa} funciona también aquí: el precursor del dolor empieza a provocar un error de antemano, y la red aprende a evitar lo que lleva al dolor. Los experimentos en personas muestran que la actividad cerebral durante el aprendizaje con dolor sigue precisamente el error de diferencias temporales~\cite{Seymour2004}, y parte de las neuronas \nt{DOPA} responde también a sucesos desagradables (capítulo~\ref{sec:dofaalt}).

\emph{Interrupción.} El dolor arranca la atención de la tarea en curso: es su función, no un efecto secundario~\cite{EcclestonCrombez1999}. En el capítulo~\ref{sec:nora} se discutió el mismo papel de “interrupción” para la ráfaga fásica de \nt{NORA}. Y la respuesta más rápida ni siquiera llega al cerebro: retirar la mano de algo caliente se cierra directamente en la médula espinal, con el mismo par de músculos y el mismo veto del antagonista de la fig.~\ref{fig:antagonists}.

\begin{keyblock}
\begin{proposition}[Señal de alarma]
\label{prop:pain}
El dolor no es una medición, sino una señal de alarma de alta prioridad. Se adapta mucho menos que los canales de medición, viaja por dos cables (el rápido dice “dónde”; el lento, “qué tan grave”), pasa por una compuerta que el tacto cierra (y que el dolor fuerte, según el supuesto de la teoría de la compuerta, abre), y su volumen lo fijan desde arriba los canales de difusión. Estos dispositivos están medidos, salvo la apertura de la compuerta por el dolor; los modelos simples~\eqref{eq:gate} y~\eqref{eq:sensitization} son simplificaciones.
\end{proposition}
\end{keyblock}

\section{Resumen}

\begin{table}[t]
\centering
\footnotesize
\setlength{\tabcolsep}{4pt}
\renewcommand{\arraystretch}{1.15}
\begin{tabular}{>{\raggedright}p{3.0cm}>{\raggedright}p{4.9cm}>{\raggedright\arraybackslash}p{4.9cm}}
\toprule
Qué hace el dolor & Cómo & Qué se sabe \\
\midrule
da la alarma & umbral alto, adaptación más débil que la del tacto & umbral medido; comparación de la adaptación: estimación \\
informa “dónde” y “qué tan grave” & dos cables de distinta velocidad~\eqref{eq:paintime} & medido \\
pasa por una compuerta & veto mediante \nt{GABA}/\nt{GLYC}~\eqref{eq:gate} & compuerta medida; apertura por el dolor: supuesto de la teoría; modelo: simplificación \\
cambia de volumen & \nt{SERO}, \nt{NORA} y opioides descendentes & medido; regla para elegir el volumen: hipótesis \\
aprende tras una lesión & aumento de la ganancia~\eqref{eq:sensitization} & sensibilización medida; modelo: simplificación \\
enseña a evitar & recompensa negativa en~\eqref{eq:ctd} & semejanza con el error de diferencias temporales medida \\
interrumpe el trabajo & captura de la atención; reflejo de retirada & medido \\
\bottomrule
\end{tabular}
\caption{El dolor como señal de alarma.}
\label{tab:pain}
\end{table}

El dolor está armado con piezas que ya vimos (tabla~\ref{tab:pain}): sensores \nt{GLUT}, dos cables, un veto mediante una intermediaria inhibidora, una perilla de volumen de difusión, un circuito RC lento, el error de predicción y la interrupción. Solo tiene una novedad: es la única entrada de la máquina cuyo volumen fija la propia máquina, una alarma que su dueño puede tanto atenuar como reconfigurar.

\section{Ejercicios}

\begin{exercise}[\lvlA]
\label{ex:14:1}
\emph{Dos cables.} La señal viene del pie, $L=1$~m. (a)~Una descarga viaja por un haz de cables de un mismo tipo cuyas velocidades llenan el rango~\eqref{eq:paintime}. ¿Cuánto se estira al llegar a la médula espinal para el tipo rápido y para el lento? (b)~Las ondas de dolor por cables de $15$ y $1$~m/s se distinguen con una diferencia desde $0.1$~s. ¿Desde qué distancia se funden?
\end{exercise}

\begin{exercise}[\lvlB]
\label{ex:14:2}
\emph{Cuando el tacto duele.} Compuerta~\eqref{eq:gate} con los parámetros del texto; la ganancia $g$ crece con la sensibilización. ¿Hasta qué $g$ el tacto sin dolor no pasa la compuerta con ninguna intensidad $\mu$? ¿Con qué $g$ empieza a pasar un tacto $\mu=1$?
\end{exercise}

\begin{exercise}[\lvlB]
\label{ex:14:3}
\emph{Estabilidad de la sensibilización.} Entradas constantes, salida de la compuerta lineal en $g$: $z=gA-C$, $A=\nu+w_{z\mu}\mu$, $C=\beta q\ge0$. (a)~Halle el equilibrio $g^*$ del modelo~\eqref{eq:sensitization} y la condición de su estabilidad. (b)~Con $k_s=0.5$, halle la intensidad de dolor por encima de la cual el modelo se dispara, sin tacto y con tacto $\mu=1$.
\end{exercise}

\begin{exercise}[\lvlA]
\label{ex:14:4}
\emph{El precursor del dolor.} Una señal en el instante $0$ predice dolor en el instante $T=2$~s; en~\eqref{eq:ctd}, $r(t)=-P\,\delta(t-T)$, $\tau_\gamma=2$~s. (a)~Halle el área del pico de $\delta$ en el instante de la señal. (b)~Una acción en el instante $t_e=1$~s cancela el dolor. ¿Qué pico de $\delta$ hay en ese instante y qué le enseñará a la red?
\end{exercise}

\begin{exercise}[\lvlB]
\label{ex:14:5}
\emph{Una alarma que se traba.} Complete la compuerta de \texttt{models/\allowbreak pain\_gate.py} con la dinámica~\eqref{eq:sensitization} y la saturación $z\le4$. El tacto $\mu=1$ está siempre presente; la lesión $\nu=2$ dura un tiempo $T$, y luego $\nu=0$. Simulando, halle el menor $T$ (en unidades de $\tau_s$) tras el cual la alarma queda encendida, con $k_s=4$, $4.6$, $5$, $6$, $8$.
\end{exercise}

\begin{exercise}[\lvlB]
\label{ex:14:6}
\emph{Diseño de la compuerta.} Con $\beta=1.5$, $w_{z\mu}=0.3$, $g=1$, elija en~\eqref{eq:gate} $q_0$, $w_{q\nu}$, $w_{q\mu}$ de modo que el umbral del dolor sea $0.2$ sin tacto y $1.0$ con tacto $\mu=1$, y que el tacto deje de atenuar el dolor en $\nu_\times=2.5$. ¿Hasta qué ganancia $g$ el tacto sin dolor no pasa la compuerta?
\end{exercise}

\begin{exercise}[\lvlC]
\label{ex:14:7}
\emph{La perilla de volumen.} El trabajo aporta un valor $V$ por unidad de tiempo, y trabajar con una lesión $\nu$ cuesta $c\nu$, así que conviene parar cuando $\nu>V/c$; el dolor interrumpe cuando $z>\theta=0.5$. (a)~¿Qué ganancia $g^*$ de la compuerta~\eqref{eq:gate} sin tacto da ese umbral con $V/c=0.25$, $0.5$, $2$? (b)~¿Con qué señales del libro podría la máquina estimar $V$ y $c$?
\end{exercise}
